\documentclass[UTF8,12pt]{ctexart}

\usepackage[a4paper,margin=1in]{geometry}
\usepackage{amsmath,amssymb,bm}
\usepackage{booktabs}
\usepackage{array}
\usepackage{hyperref}
\usepackage{enumitem}

\hypersetup{
    colorlinks=true,
    linkcolor=black,
    urlcolor=blue,
    citecolor=black
}

\setlist[itemize]{leftmargin=2em,itemsep=0.2em,topsep=0.2em}
\setlist[enumerate]{leftmargin=2em,itemsep=0.2em,topsep=0.2em}

\title{期中项目}
\author{张研疏 \quad 2200010823}
\date{\today}

\begin{document}

\maketitle

\section{数据访问、清洗与参数估计}

本节对应项目要求中的数据访问、数据清洗与输入参数估计部分。本文使用真实市场日频价格数据构造资产收益率矩阵，并在训练窗口上估计组合优化模型所需的期望收益向量与协方差矩阵。设资产数量为 $n$，清洗后的交易日数量为 $T+1$，第 $i$ 个资产在第 $t$ 个交易日的复权价格记为 $p_{t,i}$。价格矩阵记为
\[
P=
\begin{pmatrix}
p_{0,1} & p_{0,2} & \cdots & p_{0,n}\\
p_{1,1} & p_{1,2} & \cdots & p_{1,n}\\
\vdots & \vdots & & \vdots\\
p_{T,1} & p_{T,2} & \cdots & p_{T,n}
\end{pmatrix}\in\mathbb{R}^{(T+1)\times n}.
\]
由价格矩阵构造收益率矩阵后，再进一步估计 $\hat\mu$ 与 $\hat\Sigma$，作为后续二次规划、二阶锥规划、ADMM 与 PDHG 算法的输入。

\subsection{数据可用性审计}

本文选择 Yahoo Finance 作为市场数据来源，并通过 Python 中的 \texttt{yfinance} 接口获取数据。数据频率为日频，资产池选取一组流动性较高的美国股票或 ETF。原始数据字段通常包括开盘价、最高价、最低价、收盘价、成交量以及复权价格等。由于本项目关注的是投资组合的收益与风险估计，实际使用的核心字段为复权收盘价。使用复权价格的原因是它能够在一定程度上修正分红、拆股等公司行为对价格序列造成的机械跳变，从而使由价格构造出的收益率更接近真实可投资收益。

在可用性方面，本文使用的接口不需要额外申请 API token，适合复现实验。数据时间跨度取为多年日频历史数据，因此可以按照严格时间顺序划分训练集、验证集和测试集，也可以进一步扩展为滚动窗口回测。项目所需的关键字段为交易日期、资产代码与复权价格，这些字段均可以通过该数据源获得。本文在下载阶段保留原始价格数据，并输出每个资产的缺失比例，以便检查资产历史长度是否足够、是否存在大面积缺失以及是否需要从资产池中剔除某些资产。

\subsection{资产池、数据清洗与收益率构造}

设初始资产池为 $\mathcal{U}_0$，其中包含 $n_0$ 个候选资产。下载得到的原始价格数据首先按照交易日期对齐，使所有资产共享同一个交易日索引。对齐后的价格矩阵中可能存在缺失值，其来源可能包括资产上市时间较晚、个别交易日缺失、停牌、数据源错误或历史记录不完整等。本文采用如下清洗原则：首先计算每个资产价格序列的缺失比例；若某个资产的缺失比例超过给定阈值，则将其从资产池中剔除；对于剩余资产中的少量缺失值，使用前向填充与后向填充进行处理；最后删除仍然存在缺失值的交易日。清洗后的资产池记为 $\mathcal{U}$，资产数量记为 $n$。

在得到清洗后的复权价格矩阵 $P$ 后，本文使用简单收益率构造收益率矩阵。对第 $i$ 个资产，第 $t$ 个收益率定义为
\[
r_{t,i}=\frac{p_{t,i}-p_{t-1,i}}{p_{t-1,i}},\qquad t=1,2,\ldots,T,\quad i=1,2,\ldots,n.
\]
因此，第 $t$ 个交易区间上的收益率向量为
\[
r_t=(r_{t,1},r_{t,2},\ldots,r_{t,n})^\top\in\mathbb{R}^n,
\]
收益率矩阵为
\[
R=
\begin{pmatrix}
r_{1,1} & r_{1,2} & \cdots & r_{1,n}\\
r_{2,1} & r_{2,2} & \cdots & r_{2,n}\\
\vdots & \vdots & & \vdots\\
r_{T,1} & r_{T,2} & \cdots & r_{T,n}
\end{pmatrix}\in\mathbb{R}^{T\times n}.
\]
这里每一行表示一个交易区间内所有资产的收益率，每一列表示一个资产的收益率时间序列。若需要降低极端观测对协方差估计的影响，可以对收益率进行分位数截尾处理，例如将每个资产收益率序列低于 $1\%$ 分位数的值替换为该分位数，高于 $99\%$ 分位数的值替换为该分位数。该处理属于数据清洗步骤，而不是优化模型本身的一部分。

\subsection{时间划分与避免信息泄露}

为了保证后续回测具有样本外意义，本文按照时间顺序划分数据，而不是随机划分。记清洗后的收益率序列为 $\{r_t\}_{t=1}^T$。本文将时间区间划分为训练集、验证集和测试集三部分：训练集用于估计期望收益向量与协方差矩阵；验证集用于选择模型参数，例如风险厌恶系数 $\gamma$、标准差上界 $\sigma_{\max}$、ADMM 罚参数 $\rho$ 或 PDHG 步长；测试集只用于最终样本外评估。

这种划分方式的关键是避免未来信息泄露。具体来说，在任意测试时点之前，只能使用该时点之前已经发生的收益率数据估计 $\hat\mu$ 和 $\hat\Sigma$，不能使用测试区间或未来区间的信息。若采用固定训练/验证/测试划分，则所有模型选择应当在训练集和验证集上完成；若采用滚动窗口回测，则每次调仓时都只使用该调仓日前的滚动训练窗口估计参数，并在下一段持有期上记录样本外收益。

\subsection{期望收益向量与协方差矩阵估计}

组合优化模型的输入参数为期望收益向量 $\hat\mu\in\mathbb{R}^n$ 与协方差矩阵 $\hat\Sigma\in\mathbb{R}^{n\times n}$。设训练窗口中共有 $T_{\mathrm{train}}$ 个收益率观测，记为 $r_1,\ldots,r_{T_{\mathrm{train}}}$。本文采用样本均值估计期望收益：
\[
\hat\mu=\frac{1}{T_{\mathrm{train}}}\sum_{t=1}^{T_{\mathrm{train}}}r_t.
\]
其中 $\hat\mu_i$ 表示第 $i$ 个资产在训练窗口中的平均日收益率。

协方差矩阵使用无偏样本协方差估计：
\[
\hat\Sigma=\frac{1}{T_{\mathrm{train}}-1}
\sum_{t=1}^{T_{\mathrm{train}}}(r_t-\hat\mu)(r_t-\hat\mu)^\top.
\]
矩阵 $\hat\Sigma$ 的对角元 $\hat\Sigma_{ii}$ 表示第 $i$ 个资产收益率的样本方差，非对角元 $\hat\Sigma_{ij}$ 表示第 $i$ 与第 $j$ 个资产收益率之间的样本协方差。由于实际数据中资产收益率可能高度相关，样本协方差矩阵可能病态甚至接近奇异。为提高数值稳定性，本文在必要时使用如下正则化协方差矩阵：
\[
\hat\Sigma_\varepsilon=\hat\Sigma+\varepsilon I,
\qquad \varepsilon>0.
\]
该处理保持矩阵对称，并将特征值整体向右平移，有助于改善后续二次规划和线性方程求解中的条件数问题。

上述估计量直接进入后续投资组合模型。对于投资组合权重 $x\in\mathbb{R}^n$，其预测收益与预测方差分别为
\[
\hat r(x)=\hat\mu^\top x,
\qquad
\hat\sigma^2(x)=x^\top\hat\Sigma x.
\]
因此，本节完成的数据处理与参数估计步骤为后续最小方差模型、二次效用模型和标准差约束模型提供了统一的数学输入。



\end{document}
